La literatura de arquitecturas \acs{saas} distingue varias estrategias para
aislar los datos de los distintos \dfnpl{tenant}, que van desde una
instancia completa de la aplicación por \dfn{tenant} hasta una única base
de datos compartida con un campo discriminador
(\texttt{organization\_id}) en cada tabla \cite{golding2024multitenant}.
El SCDEE adopta esta última, motivada por la simplicidad operativa y la
eficiencia de recursos: una única instancia de cada servicio atiende a
todas las organizaciones, y el aislamiento se garantiza a nivel lógico
mediante el filtrado automático por \texttt{organization\_id} que aplica
el \acs{orm} en todas las consultas. Esta decisión reconoce su compromiso
(menor aislamiento físico que en estrategias más segregadas) y deja
abierta la migración a aislamientos más fuertes (por ejemplo, esquemas
PostgreSQL separados por \dfn{tenant}) como línea de trabajo futuro sin
cambios en el código de negocio.
